\scanpage{19}\mspage{17}
D'ailleurs, les deux assertions \inserted{(sur $f',g'$, sans préciser $X,\unsure{x|a},y,z$, et sur $\varphi,\psi,\alpha$)} sont en fait équivalentes,
comme on voit en faisant $X=\unsure{\mathscr{L}|\mathcal{L}}_{m+n},\ x=a,y=c,z=b$.
Ceci \unsure{se trouve|est contenu} dans le lemme \deleted{\unsure{préc.}} plus général

\uline{Lemme}. Soient $k$ et $l$ deux entiers. Alors la
sous-catégorie pleine de $\underline{\mathrm{Hom}}(\mathscr{L}_k,\mathscr{L}_l)$ formée des
foncteurs \deleted{\unsure{qui}} $\mathscr{L}_k\to\mathscr{L}_l$ appliquant source en source,
but en but est connexe. (NB Elle est non vide
\unsure{ssi} $k\geq l$).
% slanted
\margin{Dans les considérations qui suivent, \unsure{on considère|on utilise|où en vérité} en fait \unsure{tacitement que} les foncteurs envisagés sont \uline{croissants} sur les sommets (pour l'ordre de leur numérotation). Donc il faudrait écrire $\underline{\mathrm{Hom}}_{\mathrm{cr}}$ au lieu de $\underline{\mathrm{Hom}}$, dans les six énoncés qui suivent, y compris dans le lemme de connexité $\unsure{h'|k'}$ \deleted{\unsure{pour}} ci-contre.}

\deleted{\unsure{Un} tel foncteur, est décrit par une suite de sous-intervalles $I_\alpha$ de $\unsure{[1,2k+1]}$
l'intervalle d'entiers $[\unsure{1|2},2l]\simeq\mathrm{Fl}(\mathscr{L}_k)$, les
$I_\alpha$ de longueur $\lambda_\alpha$ paire, \inserted{mutuellement disjoints,} et avec $\sum\lambda_\alpha+l=k$,
et tel que le plus petit élément de $[1,2l]$ \unsure{compris dans} \inserted{qui ne soit dans aucun des}
$I_\alpha$ soit une flèche de même sens que la première flèche de $\mathscr{L}_l$.
Pour}
\[
\deleted{(I_\alpha)\in\unsure{\mathfrak{P}|\mathcal{P}}(\unsure{\mathfrak{P}|\mathcal{P}}(\mathscr{L}_k))\quad\unsure{;;}}
\]

On va décrire entièrement ces foncteurs \inserted{(\unsure{strt.|tacit.} croissants)}
$\mathscr{L}_k\xrightarrow{\varphi}\mathscr{L}_l$, en identifiant $\mathrm{Fl}(\mathscr{L}_k)$ \unsure{;|à} l'intervalle
d'entiers $[1,2k]$, et en associant à $\varphi$ la partie
\[
E_\varphi\subset\mathrm{Fl}(\mathscr{L}_k)\simeq[1,2k]
\]
\margin{$E_\varphi=\{u\in\mathrm{Fl}(\unsure{\mathscr{L}}_k)\mid\varphi(u)\ \text{une identité}\}$}
formée des $u\in\mathrm{Fl}(\mathscr{L}_k)$ \deleted{telles que} qui sont ``contractées''
par $\varphi$ en une identité. Si $E$ est un ensemble
d'entiers, on y introduit la notion de composantes
connexes (par la relation \inserted{d'équiv. de $E$} engendrée par la relation
de succession $y=x+1$). On voit aussitôt que
les composantes connexes de $E_\varphi$ \inserted{qui ne contiennent ni la première ni la dernière flèche de $\mathscr{L}_k$} sont paires, et
que
\[
\underbrace{\operatorname{card}\unsure{\mathrm{Fl}|\mathrm{El}}(\mathscr{L}_k)}_{2k}
=\operatorname{card}\lvert E_\varphi\rvert+\underbrace{\operatorname{card}\mathrm{Fl}(\mathscr{L}_{\unsure{l|a}})}_{2l}
\]
% archive sheet number: 18
